% =============================================
% Section 5: Xtensa-Native Bitwise Inference Kernel
% =============================================
\section{Xtensa-Native Bitwise Inference Kernel}
\label{sec:kernel}

This section describes the design and implementation of our Xtensa LX7-native
ternary inference kernel, optimized for the \espss{}'s SIMD instruction set.

\subsection{Xtensa LX7 Architecture Overview}

The Xtensa LX7 processor in the \espss{} is a 32-bit RISC core with
configurable extensions. Espressif's implementation includes the EE (ESP
Extension) SIMD instruction set, which provides 128-bit vector operations:

\begin{itemize}[leftmargin=*,topsep=2pt,itemsep=1pt]
  \item \texttt{EE.VMULAS.S8.QACC}: 8-bit signed multiply-accumulate
  \item \texttt{EE.VZIP.8}: 8-bit vector interleave
  \item \texttt{EE.VCMP.GT.S8}: 8-bit signed vector comparison
  \item \texttt{EE.MOVI.32.A}: 128-bit vector load from aligned memory
\end{itemize}

Critically, the Xtensa LX7 lacks a hardware FPU---all floating-point
operations are software-emulated, making them approximately 10--20$\times$
slower than integer operations. This hardware constraint makes ternary
(integer-only) inference not merely a compression technique but a
\emph{fundamental requirement} for practical performance.

\subsection{Ternary GEMV Kernel Design}

The core operation in Transformer inference is the General Matrix-Vector
multiplication (\gemv{}): $\mathbf{y} = W\mathbf{x}$, where
$W \in \{-1, 0, +1\}^{m \times n}$ and $\mathbf{x} \in \mathbb{Z}^n$
(quantized activations).

\subsubsection{Weight Packing}

We encode ternary weights using 2 bits per weight, as defined in
Equation~\ref{eq:ternary}. For a weight matrix of dimensions $m \times n$,
this requires $\frac{2mn}{8} = \frac{mn}{4}$ bytes of storage, representing
a $16\times$ compression over FP32.

For efficient SIMD processing, weights are packed into 128-bit vectors,
each containing 64 ternary weights:

\begin{equation}
\text{vec}_{128} = [w_0^{b_1b_0}, w_1^{b_1b_0}, \ldots, w_{63}^{b_1b_0}]
\end{equation}

\subsubsection{Bitwise GEMV Algorithm}

The ternary \gemv{} is decomposed as described in
Equation~\ref{eq:bitwise_gemv}. We implement this using three bitwise
operations per weight vector:

\begin{algorithm}[H]
\caption{Ternary \gemv{} --- single output element $y_i$}
\label{alg:ternary_gemv}
\begin{algorithmic}[1]
\REQUIRE Packed weight row $W_i$ (2-bit encoded), input $\mathbf{x}$ (INT8)
\ENSURE Output $y_i$ (INT32 accumulator)
\STATE $\text{sign\_bits} \leftarrow W_i[\text{bit}_1]$
\COMMENT{Extract sign plane}
\STATE $\text{mask\_bits} \leftarrow W_i[\text{bit}_0] \,|\, W_i[\text{bit}_1]$
\COMMENT{Non-zero mask}
\STATE $\text{pos\_sum} \leftarrow \text{masked\_sum}(\mathbf{x}, \text{mask\_bits} \,\&\, \overline{\text{sign\_bits}})$
\STATE $\text{neg\_sum} \leftarrow \text{masked\_sum}(\mathbf{x}, \text{mask\_bits} \,\&\, \text{sign\_bits})$
\STATE $y_i \leftarrow \text{pos\_sum} - \text{neg\_sum}$
\end{algorithmic}
\end{algorithm}

\subsubsection{SIMD Optimization}

Using the EE instruction set, we process 16 INT8 activations per cycle
via the 128-bit vector unit. The masked summation in
Algorithm~\ref{alg:ternary_gemv} maps to:

\begin{enumerate}[leftmargin=*,topsep=2pt,itemsep=1pt]
  \item \texttt{EE.MOVI.32.A}: Load 16 activations into vector register
  \item Bitwise AND with weight mask (unpacked to byte mask via LUT)
  \item \texttt{EE.VMULAS.S8.QACC}: Accumulate masked values
\end{enumerate}

\subsection{Memory Layout Optimization}

\subsubsection{Cache-Friendly Access Pattern}

The \espss{}'s \psram{} operates over an SPI bus with high latency for
random access but reasonable throughput for sequential access. We optimize
the memory layout for row-major sequential scanning:

\begin{itemize}[leftmargin=*,topsep=2pt,itemsep=1pt]
  \item Weight matrices are stored in row-major packed format
  \item Input vectors are aligned to 16-byte boundaries
  \item Accumulator arrays reside in internal SRAM for fast access
\end{itemize}

\subsubsection{Tiled Computation}

For large matrices that exceed the L1 cache, we employ tiling with
tile sizes chosen to maximize \psram{} burst read efficiency:

\begin{equation}
T_{\text{opt}} = \underset{T}{\arg\min} \; \frac{m \cdot n}{T} \cdot L_{\text{overhead}}
+ T \cdot L_{\text{compute}}
\end{equation}

\noindent where $L_{\text{overhead}}$ is the per-tile setup latency and
$L_{\text{compute}}$ is the per-element computation latency.

\subsection{Implementation Summary}

    Our benchmark results on the \espss{} show that the bitwise \gemv{} kernel
    provides a \textbf{3.2$\times$} speedup and reduces PSRAM bandwidth consumption
    by 75\% compared to standard floating-point implementations. This efficiency
    is direct evidence of the advantage of mapping ternary weights to 
    SIMD-friendly bitwise primitives.
